\section{Shrinking two hereditary supports}
\label{shr:section}

We will shrink two hereditary supports so that a prescribed
cross-compression decreases by a fixed factor, while at least one new
support remains nonzero in each irreducible representation where an
original element has norm at least $1/2$.  Repeated use of this operation
will separate the supports
in the norm-cone construction of Section~\ref{norm:section}.
Only the two Cuntz-semigroup properties in \textup{(C)} are needed;
nuclearity is not used here.  All
Cuntz comparisons in hereditary subalgebras are understood in the ambient
stabilization, with their witnesses compressed into the indicated corners
as in the preliminaries.  If $\sigma$ is a representation, its matrix
amplifications will also be denoted by $\sigma$.

\begin{lemma}[Two-support shrinking]\label{lem:shrinking}
There is a numerical constant $0<c<1$ with the following property for
every $C^*$-algebra $A$ satisfying \textup{(C)}.  If $a,b\in A$ are positive
contractions and $z\in\widetilde A$ satisfies
$\|p_a z p_b\|\le1$, there are positive contractions
$a'\in\her(a)$ and $b'\in\her(b)$ such that
\begin{equation}\label{shr:conclusion}
 \|p_{a'}z p_{b'}\|\le c
\end{equation}
and, for every irreducible representation $\sigma$ of $A$,
\begin{equation}\label{shr:visibility}
 \max\{\|\sigma(a)\|,\|\sigma(b)\|\}\ge\tfrac12
 \quad\Longrightarrow\quad
 \sigma(a')\ne0\ \text{or}\ \sigma(b')\ne0.
\end{equation}
The support norms are computed in the bidual, or equivalently in the
universal representation.
\end{lemma}

We first prove a transfer statement for two hereditary supports.  It
replaces finitely many rectangular elements by one contraction which
has norm one in every representation where a prescribed element reaches
a fixed norm threshold.

\begin{lemma}[Transfer between two supports]\label{lem:transfer}
Suppose that $A$ satisfies \textup{(C)}.  Let $B\subseteq A$ be hereditary,
let $D\in M_l(B)_+$ and $E\in B_+$, and suppose that
$w_1,\ldots,w_m\in M_{l,1}(B)$ satisfy
\[
 w_iw_i^*\in\her(D),\qquad w_i^*w_i\in\her(E).
\]
For each $\delta>0$ there is a contraction $v\in M_{l,1}(B)$ such that
\[
 vv^*\in\her(D),\qquad v^*v\in\her(E),
\]
and, in every representation $\sigma$ of $B$,
\[
 \max_{1\le i\le m}\|\sigma(w_i)\|\ge\delta
 \quad\Longrightarrow\quad \|\sigma(v)\|=1.
\]
\end{lemma}

\begin{proof}
Set $T_0=\operatorname{diag}(w_1^*w_1,\ldots,w_m^*w_m)$.  The hereditary
conditions and $[w_i^*w_i]=[w_iw_i^*]$ give
\[
 [T_0]\le m[E],\qquad [T_0]\le m[D].
\]
By divisibility in \textup{(C)}, choose a positive element $r$ in a
stabilization with
\[
 (m+1)[r]\le[T_0]\le(m+2)[r].
\]
Almost unperforation, applied with $m$, gives $[r]\le[E]$ and
$[r]\le[D]$.  Choose $\eta>0$ sufficiently small that
\begin{equation}\label{shr:transfer-cut}
 [(T_0-\delta^2/2)_+]\le(m+2)[(r-2\eta)_+].
\end{equation}
This is possible because the class on the left is compactly contained in
$[T_0]$, and the classes $(m+2)[(r-2\eta)_+]$ increase to $(m+2)[r]$.

The cut-placement result in the preliminaries supplies partial isometries
$u_E,u_D$ which identify $\her((r-\eta)_+)$ with hereditary subalgebras of
$\her(E)$ and $\her(D)$, respectively.  Their initial projection is the
support of $(r-\eta)_+$.  Choose a continuous function $j$ with
$0\le j\le1$, $j=0$ on $[0,\eta]$, and $j=1$ on $[2\eta,\infty)$.
Then
\[
 v=u_D j(r)u_E^*
\]
belongs to $M_{l,1}(B)$: factor it as
$(u_Dj(r)^{1/2})(u_Ej(r)^{1/2})^*$.  Each factor belongs to the ambient
stabilized algebra by cut placement, and the product's left and right
supports lie in the prescribed finite matrix corners.  It is a contraction, and
\[
 v^*v=u_Ej(r)^2u_E^*\in\her(E),\qquad
 vv^*=u_Dj(r)^2u_D^*\in\her(D).
\]
If $\|\sigma(w_i)\|\ge\delta$ for some $i$, then
$\sigma((T_0-\delta^2/2)_+)\ne0$.  By
\eqref{shr:transfer-cut}, the source copy
$u_E(r-2\eta)_+u_E^*$ is nonzero in this representation: otherwise the
comparison would make the represented left side zero.  The operator
$\sigma(v^*v)$ is the identity on the support of that nonzero copy, since
$j$ is one on its spectrum.  Thus $\|\sigma(v)\|=1$.
\end{proof}

\begin{proof}[Proof of Lemma~\ref{lem:shrinking}]
Suppose that no such constant exists.  Fix numbers $c_n\in(0,1)$ increasing
to $1$.  For each $n$, choose an algebra satisfying \textup{(C)} and data
$a,b,z$ for which the conclusion fails with $c=c_n$.  The algebras are
allowed to vary; we suppress $n$ until the limiting argument.  The
construction below always gives \eqref{shr:conclusion}, so failure will
provide a representation in which both constructed positives vanish.

\medskip
\noindent\textbf{Step 1: a matrix corner and two prospective supports.}
Apply Lemma~\ref{lem:placement}, with seven slots and a cut of $a$ below
$1/2$, to obtain a c.p.c.\ order-zero map $M_7\to\her(a)$ with weight $d$
satisfying
\begin{equation}\label{shr:d-visibility}
 \|\sigma(a)\|\ge\tfrac12\quad\Longrightarrow\quad\|\sigma(d)\|=1.
\end{equation}
Its supporting matrix units identify
\[
 \her(d)=M_7(B),\qquad d=\operatorname{diag}(P),
\]
where $B$ is the hereditary algebra generated by its first diagonal
weight and $P\in B$ is a strictly positive contraction.  Concretely, in
any representation, the support space of $d$ is identified with
$\mathcal H^{\oplus7}$, where $B$ acts nondegenerately on $\mathcal H$.
Choose $Q=\min(1,2b)$ by functional calculus.  Then $Q\in\her(b)$ is a
positive contraction and
\begin{equation}\label{shr:Q-visibility}
 \|\sigma(b)\|\ge\tfrac12\quad\Longrightarrow\quad\|\sigma(Q)\|=1.
\end{equation}

We shall choose contractions
$v_0\in M_{2,1}(B)$ and $v_1\in M_{3,1}(B)$.  For any such choice set,
in multiplier notation,
\begin{equation}\label{shr:U}
 U=\begin{pmatrix}
 \frac12\,1\\[2pt] \frac12 v_0\\[2pt] \frac12 v_1\\[2pt]
 \bigl(\frac34\,1-\frac14v_0^*v_0-\frac14v_1^*v_1\bigr)^{1/2}
 \end{pmatrix},\qquad
 L=UPU^*\in M_7(B),\qquad t=\tfrac12.
\end{equation}
The expression under the square root is at least $\frac14\,1$, and
$U^*U=1$.  Define
\begin{equation}\label{shr:ab-prime}
 a'=(L-t)_+,\qquad
 b'=Q\bigl[Q(c_n^2t-z^*Lz)Q\bigr]_+Q.
\end{equation}
These belong to $\her(a)$ and $\her(b)$, respectively.  They are
contractions: $0\le L\le1$, and the selfadjoint operator inside the positive
part in the second expression is bounded above by $c_n^2t\,1$.
By Lemma~\ref{lem:support}, the compression of $z^*Lz$ to the support of
$b'$ is at most $c_n^2t$.  Since $L\ge t p_{a'}$, this implies
\[
 t p_{b'}z^*p_{a'}z p_{b'}
 \le p_{b'}z^*Lz p_{b'}\le c_n^2t p_{b'},
\]
which proves \eqref{shr:conclusion} with $c_n$, for every choice of the two
columns.

\medskip
\noindent\textbf{Step 2: data chosen before either column.}
Put $Z=p_d zQ$ and write its rows in groups of sizes $1,2,3,1$ as
$Z_{\mathrm f},Z_0,Z_1,Z_{\mathrm{last}}$.  These are operators from the
original representation space to the corresponding copies of
$\mathcal H$.  The hypothesis on $z$ gives
\begin{equation}\label{shr:Z-bound}
 Z^*Z\le Q^2,
\end{equation}
because $p_d\le p_a$ and $p_bQ=Q$.  Thus $Z$ has norm at most one, even
though no uniform bound on $z$ itself is assumed.

The next choices address the only way the prospective supports can fail
to retain visibility.  For any choice of the columns, a representation
in which $a'=b'=0$ would have $P\le t$ and, by the zero case of
Lemma~\ref{lem:support},
\begin{equation}\label{shr:alignment-chain}
 c_n^2tQ^2\le Z^*UPU^*Z\le tZ^*Z\le tQ^2.
\end{equation}
As $c_n\to1$, these inequalities force the limiting range of $Z$ into
the range of $U$, with its $U$-coordinate in the $t$-spectral subspace
of $P$.  Step~5 proves this alignment precisely.  Since the first row
of $U$ is $1/2$, the first row of $Z$ then recovers that coordinate.
The following cutoffs implement this recovery within the hereditary
corner, before either column has been chosen.

Fix continuous functions on $[0,1]$ by
\[
 g(x)=\min\{1,\max\{0,8x-2\}\},\qquad G(x)=\min\{1,4x\}.
\]
They vanish at zero, satisfy $Gg=g$, and $g(t)=G(t)=1$.  Set
\begin{equation}\label{shr:XY}
 X=2g(P)Z_{\mathrm f},\qquad h_0=XX^*,\qquad
 Y_i=2\operatorname{diag}(G(P))Z_iX^*\quad(i=0,1).
\end{equation}
Here $h_0\in B_+$, $Y_0\in M_{2,1}(B)$, and
$Y_1\in M_{3,1}(B)$.  For example, their entries are the appropriate
entries of sandwiches of $zQ^2z^*$ by $g(d)$ and $G(d)$; these sandwiches
belong to $\her(d)$.  We have
\begin{equation}\label{shr:initial-bounds}
 \|X\|\le2,\quad \|h_0\|\le4,\quad \|Y_i\|\le4,
 \qquad G(P)h_0=h_0=h_0G(P).
\end{equation}
In particular, $G(P)$ acts as the identity on the support of $h_0$.
None of these data depends on $v_0$ or $v_1$.

For later use fix
\[
 r=\tfrac18,\qquad
 W_j=r^{-1}(Y_0)_j^*\quad(j=1,2),\qquad
 f=f(h_0),\quad f(x)=\min\{1,\max\{0,4x-1\}\}\quad(x\ge0).
\]
Thus $f$ vanishes on $[0,1/4]$ and is one on $[1/2,\infty)$.
With $u=W_1f$ and $c_f=(1-f^2)^{1/2}$, put $V=u+c_f$ in
$\widetilde B$ and define
\begin{equation}\label{shr:T}
 T=1-VV^*=f^2-uu^*-u c_f-c_f u^*\in B.
\end{equation}
This selfadjoint element records the prospective range complement of
$V$.  At this stage $V$ need not be an isometry and $T$ need not be
positive.  Step~7 will produce a reducing subspace in the failure limit
on which $V$ is an isometry and $T$ is a nonzero projection.  We use
the positive spectral part of $T$ now to prepare the second column,
which will rule out that possibility.
The bounds above give
$\|W_j\|\le32$, $\|u\|\le32$, and $\|T\|\le1089$.  Fix on $\R$ the
continuous functions
\[
 \beta(x)=\min\{1,\max\{0,2x-1\}\},\qquad
 \chi(x)=\min\{1,\max\{0,4x-1\}\}.
\]
They vanish at zero, $\beta$ vanishes on $(-\infty,1/2]$,
$\beta(1)=1$, and $\chi\beta=\beta$.  Put
\begin{equation}\label{shr:BT-y}
 b_T=\beta(T),\qquad B_T=\her_B(b_T),\qquad
 y=\operatorname{diag}(b_T)Y_1b_T\in M_{3,1}(B_T),\qquad
 S=\operatorname{diag}(b_T)\in M_3(B_T).
\end{equation}
We have $\|y\|\le4$, and $S$ is strictly positive in $M_3(B_T)$.

\medskip
\noindent\textbf{Step 3: the second column and bounded comparison witnesses.}
Apply Lemma~\ref{lem:placement} with two slots inside
$\her((yy^*-10^{-3})_+)$, using the narrower cut
$(yy^*-10^{-2})_+$.  The plateau of the slot weights may be chosen equal
to one on the transported narrower element.  To use the contraction
form of that Lemma explicitly, take
$b=(yy^*-10^{-3})_+/16$ and $\epsilon=9\cdot10^{-3}/16$;
positive rescaling preserves all the hereditary algebras and Cuntz
classes in question.  We obtain equivalent
orthogonal positive contractions $k,k'\in M_3(B_T)$ with
\begin{equation}\label{shr:k-comparison}
 k,k'\in\her((yy^*-10^{-3})_+),\qquad
 [(yy^*-10^{-2})_+]\le3[(k-0.9)_+].
\end{equation}
If this hereditary algebra is zero, take $k=k'=0$; the same assertions
hold.  Choose three matrices $x_j\in M_3(B_T)$ witnessing
\[
 \left\|(yy^*-0.01)_+-\sum_{j=1}^3x_j^*(k-0.9)_+x_j\right\|<0.005.
\]
There are three matrices because the right side of the comparison is a
threefold direct sum; rectangular witnesses can be split into its three
blocks.  Replace $x_j$ by
\[
 q_j=\left(\frac{(k-0.9)_+}{(k-0.8)_+}\right)^{1/2}x_j,
\]
where the quotient is defined to be zero when its denominator is zero.
This is a continuous function of $k$.  Then
\begin{equation}\label{shr:bounded-witnesses}
 \left\|(yy^*-0.01)_+-\sum_{j=1}^3q_j^*(k-0.8)_+q_j\right\|<0.005,
 \qquad \|q_j\|\le\sqrt{160}<13.
\end{equation}
Indeed the compression sum before replacement is positive and has norm
less than $16$, so each summand has norm less than $16$.  Also
\[
 \frac{(k-0.9)_+}{(k-0.8)_+}\le10(k-0.9)_+.
\]
These give $q_j^*q_j\le10x_j^*(k-0.9)_+x_j$ and the stated bound.  The
positive gap between the two cuts is what makes these witnesses uniformly
bounded.

Choose $p'=\min(1,k/0.8)$ and put
\[
 D'=(1-p')S(1-p')\in M_3(B_T)_+.
\]
We claim that
\begin{equation}\label{shr:Dprime-comparison}
 3[b_T]=[S]\le[k]+[D']\le2[D'].
\end{equation}
For the first inequality, sandwich
$\frac12\,1\le(p')^2+(1-p')^2$ by $S^{1/2}$.  The resulting two summands
have the classes of $p'Sp'$ and $(1-p')S(1-p')$, and
$p'Sp'\in\her(k)$.  For the second inequality it suffices to check that
$k'\in\her(D')$, since $[k]=[k']$.  Orthogonality gives
$(1-p')k'(1-p')=k'$.  Since $S$ is strictly positive, there are positive
elements $c_j\to k'$ with $c_j\le C_jS$ for finite constants $C_j$:
one may take $c_j=e_jk'e_j$ with
$e_j=S(S+1/j)^{-1}$.  Consequently
\[
 (1-p')c_j(1-p')\le C_jD',\qquad
 (1-p')c_j(1-p')\longrightarrow k',
\]
which proves the inclusion and \eqref{shr:Dprime-comparison}.

Almost unperforation now gives $[b_T]\le[D']$.  Cut placement supplies a
column $v_1\in M_{3,1}(B_T)$ with
\begin{equation}\label{shr:v1}
 v_1^*v_1=(b_T-1/4)_+,\qquad
 v_1v_1^*\in\her(D'),\qquad \|v_1\|\le1.
\end{equation}
Since $(k-0.8)_+(1-p')=0$, its range condition implies
\begin{equation}\label{shr:v1-annihilation}
 (k-0.8)_+v_1=0.
\end{equation}
We shall also retain the following fixed continuous support identities:
\begin{equation}\label{shr:fixed-supports}
 (10^{-3}-yy^*)_+k=0,\qquad
 \chi(T)(v_1)_j=(v_1)_j=(v_1)_j\chi(T)\quad(1\le j\le3).
\end{equation}
The first follows from \eqref{shr:k-comparison}; the second holds because
$\chi(T)b_T=b_T$ and the entries of $v_1$ belong to $B_T$.  These equations,
unlike statements about varying support projections, pass directly to an
operator ultraproduct.

\medskip
\noindent\textbf{Step 4: finite word tests and the first column.}
Let $p$ be the continuous function of $Y_0Y_0^*$ which is zero up to
$r^2$, linear between $r^2$ and $(r+1/n)^2$, and one thereafter.  Put
\begin{equation}\label{shr:DE}
 s=\operatorname{diag}((h_0-1/4)_+)\in M_2(B),\qquad
 E=(h_0-1/2)_+\in B,\qquad
 D=s\bigl[s(\tfrac14-p)s\bigr]_+s\in M_2(B).
\end{equation}
Choose an ordered alphabet consisting of all entries of
\begin{equation}\label{shr:alphabet}
 P,h_0,Y_0,Y_1,p,s,E,D,T,b_T,y,k,k',D',v_1,q_1,q_2,q_3
\end{equation}
and their adjoints.  The number of letters is independent of $n$, as all
matrix sizes are fixed.  Their norms are bounded independently of the
instance: the bounds already given apply, $\|s\|,\|E\|\le4$,
$\|D\|\le256$, and $p,k,k',b_T,D',v_1$ are contractions.  Thus, for
example, $1089$ is a common bound on every letter.  Neither $v_0$ nor $U$
is a letter.

Apply Lemma~\ref{lem:transfer} to the finite list of columns
\begin{equation}\label{shr:word-tests}
 D\mathbf e_jxE,\qquad j=1,2,\quad
 x\text{ a word of length at most }n,
\end{equation}
including the empty word $x=1$, with threshold $1/n$.  Here
$\mathbf e_j$ is the $j$th scalar basis column.  The left and right squared
moduli of these columns lie in $\her(D)$ and $\her(E)$, respectively.
The resulting contraction $v_0\in M_{2,1}(B)$ satisfies
\begin{equation}\label{shr:v0}
 v_0v_0^*\in\her(D),\qquad v_0^*v_0\in\her(E),\qquad
 v_0^*pv_0\le\tfrac14v_0^*v_0.
\end{equation}
The last assertion is Lemma~\ref{lem:support} applied to the pencil
defining $D$.  In addition, we have the exact operator identities
\begin{equation}\label{shr:v0-fixed}
 (1/2-h_0)_+v_0^*=0,\qquad
 \operatorname{diag}(G(P))v_0=v_0.
\end{equation}
For the first use the right hereditary condition.  For the second, the
support of $D$ is contained in that of $s$, and $G(P)$ fixes the support
of $h_0$ by \eqref{shr:initial-bounds}.  The choice of $v_1$, all the
comparison witnesses, and the whole alphabet precedes the choice of
$v_0$, so the construction is not circular.

\medskip
\noindent\textbf{Step 5: failure forces alignment in a Hilbert ultraproduct.}
Make the choices above at coordinate $n$ and define $a',b'$ by
\eqref{shr:ab-prime}.  Since they satisfy the required angle bound, the
assumed failure gives an irreducible representation $\sigma_n$ such that
\[
 \max\{\|\sigma_n(a)\|,\|\sigma_n(b)\|\}\ge\tfrac12,
 \qquad \sigma_n(a')=\sigma_n(b')=0.
\]
Use the same symbols for the represented operators.  The first vanishing
gives $L\le t$, hence $P=U^*LU\le t$.  By
\eqref{shr:d-visibility}, the $a$-trigger is impossible.  Therefore
$\|Q\|=1$ by \eqref{shr:Q-visibility}.  The two vanishings also give
the conditional inequalities \eqref{shr:alignment-chain} established
in Step~2.

Pass along the fixed free ultrafilter to Hilbert-space ultraproducts of
the representation spaces and of their first matrix-coordinate spaces.
For clarity, such an ultraproduct is the completion of bounded sequences
of vectors modulo those whose norms tend to zero along the ultrafilter,
with inner product the ultralimit of the coordinate inner products.
Uniformly bounded operator sequences act on these spaces, respect sums,
products and adjoints, and have norm equal to the ultralimit of their
norms.  The last assertion follows by choosing a nearly norm-maximizing
unit vector at each coordinate.  Fixed continuous functional calculus is
preserved, by uniform polynomial approximation on a common spectral
interval.  All operators used here are uniformly bounded: in addition to
the alphabet this includes $Q,Z,X,U,v_0$, by
\eqref{shr:Z-bound}, \eqref{shr:initial-bounds}, and \eqref{shr:U}.
We keep the same notation for their limits.

Equality holds throughout \eqref{shr:alignment-chain}.  In particular
$Z^*Z=Q^2$ and
$Z^*(t1-UPU^*)Z=0$.  Since
\[
 t1-UPU^*=t(1-UU^*)+U(t-P)U^*
\]
is a sum of positive operators, both positive summands annihilate $Z$.
For $\xi=U^*Z$, it follows that
\begin{equation}\label{shr:alignment}
 Z=U\xi,\qquad P\xi=t\xi,\qquad \xi^*\xi=Q^2.
\end{equation}
The first row of $U$ is $1/2$, so $X=2g(P)Z_{\mathrm f}=\xi$.
Consequently
\begin{equation}\label{shr:limit-Y}
 \|h_0\|=1,\qquad
 Y_i=\operatorname{diag}(G(P))v_i h_0,\qquad
 Y_0=v_0h_0.
\end{equation}
The last equality uses the second identity in \eqref{shr:v0-fixed}.

\medskip
\noindent\textbf{Step 6: the moving cutoff makes every word test vanish.}
Although the function defining $p$ changes with $n$, $p$ itself is a
bounded alphabet variable.  At each coordinate it commutes with
$R_0=Y_0Y_0^*$ and satisfies
\[
 (r^2-R_0)_+p=0,\qquad
 \|(1-p)(R_0-r^2)_+\|\le(r+1/n)^2-r^2.
\]
Thus in the limit
\begin{equation}\label{shr:moving-cutoff}
 [p,R_0]=0,\qquad (r^2-R_0)_+p=0,\qquad
 (1-p)(R_0-r^2)_+=0.
\end{equation}
The first two equations show that the support of $p$ lies in the spectral
subspace of $R_0$ for $[r^2,\infty)$, hence in the closure of the range of
$Y_0$, and therefore of $v_0$ by \eqref{shr:limit-Y}.  On this last closed
range the inequality $v_0^*pv_0\le\frac14v_0^*v_0$ says that the quadratic
form of $p$ is at most one quarter of the squared norm.  This follows
first on vectors $v_0\eta$, then by density; on the orthogonal complement
$p$ is zero.  Hence
\begin{equation}\label{shr:p-small}
 0\le p\le\tfrac14\,1.
\end{equation}
The last equation in \eqref{shr:moving-cutoff}, and invertibility of
$1-p$, now give $R_0\le r^2$.  The other support equation then gives
\begin{equation}\label{shr:R0}
 \|Y_0\|\le r,\qquad R_0p=r^2p.
\end{equation}

Every fixed word $x$ in the limiting alphabet satisfies
\begin{equation}\label{shr:all-word-tests}
 D\mathbf e_jxE=0\qquad(j=1,2).
\end{equation}
Indeed, if such an operator had norm $\delta>0$, its coordinate norms
would be at least $\delta/2$ on an ultrafilter-large set.  On an
ultrafilter-large subset we also have $n\ge\operatorname{length}(x)$ and
$1/n\le\delta/2$.  The transfer condition used in
\eqref{shr:word-tests} would give $\|v_0\|=1$ in the ultraproduct.
The first equation in \eqref{shr:v0-fixed} puts the range of $v_0^*$ in
the $h_0\ge1/2$ spectral subspace.  Inverting $h_0$ on that subspace,
with inverse norm at most two, gives
\[
 \|v_0h_0\|\ge\tfrac12\|v_0\|=\tfrac12.
\]
This contradicts $Y_0=v_0h_0$ and \eqref{shr:R0}, since $r=1/8$.

\medskip
\noindent\textbf{Step 7: a nonzero defect projection on the word space.}
Let $\mathcal H$ denote the limiting single-coordinate Hilbert space and
set
\[
 \mathcal L=\overline{\operatorname{span}}
 \{xE\eta:x\text{ a word in the alphabet},\ \eta\in\mathcal H\}.
\]
The alphabet includes adjoints, so $\mathcal L$ reduces every letter and
the $C^*$-algebra they generate.  Restrict letters to this space, and
matrix variables to its amplifications.  Equation
\eqref{shr:all-word-tests} gives $D=0$ on $\mathcal L^{\oplus2}$.
Moreover $E\mathcal H\subseteq\mathcal L$, so the restriction retains
$\|E\|=1/2$, hence $\|h_0\|=1$.  We do not assert that $v_0$ reduces
$\mathcal L$.  All further calculations involve only alphabet operators
and their continuous functions.  The needed range information from
$v_0$ has already become the identity
\begin{equation}\label{shr:W-range}
 (1/2-h_0)_+W_j=0\qquad(j=1,2),
\end{equation}
which follows on the full limiting space from
$W_j=r^{-1}h_0(v_0)_j^*$ and \eqref{shr:v0-fixed}, and thus restricts to
$\mathcal L$.

By the zero case of Lemma~\ref{lem:support}, $D=0$ implies
$s(\frac14-p)s\le0$.  Equation \eqref{shr:p-small} gives the reverse
inequality, so $s(\frac14-p)s=0$.  Taking a positive square root shows
$(\frac14-p)s=0$, or $ps=s/4$.  Together with \eqref{shr:R0}, this yields
\begin{equation}\label{shr:Rs}
 R_0s=r^2s.
\end{equation}
Let $e=1_{(1/4,\infty)}(h_0)$ on $\mathcal L$.  The range of $s$ is dense
in $(e\mathcal L)^{\oplus2}$, so \eqref{shr:Rs} implies
\[
 (W_i e)^*(W_j e)=\delta_{ij}e\qquad(1\le i,j\le2).
\]
Thus $W_1,W_2$ restricted to $e\mathcal L$ are isometries with orthogonal
ranges.  By \eqref{shr:W-range}, their ranges lie in the
$h_0\ge1/2$ subspace, where $f=1$.

It follows that $u=W_1f$ satisfies $u^*u=f^2$ and $c_fu=0$, hence also
$u^*c_f=0$.  Consequently
\[
 V^*V=f^2+c_f^2=1.
\]
Its range-complement projection is exactly
\begin{equation}\label{shr:T-projection}
 1-VV^*=f^2-uu^*-uc_f-c_fu^*=T.
\end{equation}
Both the initial and final supports of $u$ are contained in $e$, and
$f=ef$.  As $c_f$ commutes with $e$, formula \eqref{shr:T-projection}
shows $T=eTe$.  In particular
\begin{equation}\label{shr:ThT}
 Th_0T\ge\tfrac14 T.
\end{equation}
This projection is nonzero.  Indeed $W_2f\ne0$ because
$\|h_0\|=\|f\|=1$ and $W_2$ is isometric on $e\mathcal L$, whereas
\[
 V^*W_2f=fW_1^*W_2f+c_fW_2f=0.
\]
Its nonzero range is therefore contained in the range of $T$.

\medskip
\noindent\textbf{Step 8: the second column gives the contradiction.}
Since $T$ is now a projection, $b_T=\beta(T)=T$ and $\chi(T)=T$.
The fixed identities \eqref{shr:fixed-supports} give
\[
 T(v_1)_j=(v_1)_j=(v_1)_jT.
\]
They hold after restriction because $v_1$ was included in the alphabet.
As $G(P)$ fixes the support of $h_0$ and $T=eTe$, it also fixes $T$;
thus $\operatorname{diag}(G(P))v_1=v_1$.  Using
\eqref{shr:BT-y}, \eqref{shr:v1}, and \eqref{shr:limit-Y}, we obtain
\begin{equation}\label{shr:final-y}
 y=v_1Th_0T,\qquad v_1^*v_1=\tfrac34 T,\qquad
 y^*y=\tfrac34(Th_0T)^2\ge\tfrac{3}{64}T.
\end{equation}
In particular, $\|yy^*\|\ge3/64$.

The three uniformly bounded witnesses in
\eqref{shr:bounded-witnesses} survive both the ultraproduct and the
restriction, giving
\[
 \left\|(yy^*-0.01)_+-\sum_{j=1}^3q_j^*(k-0.8)_+q_j\right\|\le0.005.
\]
Hence $q=(k-0.8)_+$ is nonzero: otherwise
$\|yy^*\|\le0.015<3/64$, contradicting \eqref{shr:final-y}.
On the other hand, \eqref{shr:v1-annihilation} and
\eqref{shr:final-y} imply $qy=0$.  The first fixed support identity in
\eqref{shr:fixed-supports} puts the range of $k$, and hence of $q$, in
the spectral subspace of $yy^*$ for $[10^{-3},\infty)$.  Therefore
\[
 0=qyy^*q\ge10^{-3}q^2,
\]
forcing $q=0$.  This contradiction proves the existence of the constant
$c$ and completes the proof.
\end{proof}
